\begin{proposition}[candidate-specific holomorphic-form vanishing]
\label{prop:workbench-holomorphic-form-vanishing}
\Verified\ For the constructed threefold,
\[
 H^0(X,\Omega_X^p)=0\qquad (p=1,2,3).
\]
\end{proposition}

\begin{proof}
Over the marked cover of the smooth part, write the fibre coordinates as
$\zeta=(\zeta_1,\zeta_2)$ and
\[
 \Pi=[Z\mid I_2],\qquad
 Z=\begin{pmatrix}6\mu&\tau\\ \beta&\mu\end{pmatrix}.
\]
Translation by $\lambda\in\Lambda$ sends
$d\zeta$ to $d\zeta+\dot\Pi\lambda\,dz$.  Compactness of each torus fibre
therefore puts every translation-invariant holomorphic form into one of the
following normal forms:
\begin{align*}
 p=1:&\quad a(z)\,dz+c(z)\cdot d\zeta,
       &c(z)\dot Z(z)&=0,\\
 p=2:&\quad dz\wedge\bigl(b(z)\cdot d\zeta\bigr),\\
 p=3:&\quad c(z)\,dz\wedge d\zeta_1\wedge d\zeta_2.
\end{align*}
Indeed, fibrewise periodic coefficients are constant.  A possible affine
part in a coefficient vanishes after restriction to the real span of the
two constant columns of $\Pi$; in degree two this also kills the putative
$d\zeta_1\wedge d\zeta_2$ term because $\dot\tau$ is nonzero on the smooth
base.

The finite fillings are obtained from smooth torus families by free cyclic
quotients, so the displayed coefficients extend holomorphically across the
two elliptic lifts.  At the cusp use the verified toric charts
$t=w_0w_1w_2$.  Substitution of
\[
 d\zeta_i=(2\pi i)^{-1}dx_i/x_i,
 \qquad ds=(2\pi i)^{-1}dt/t
\]
and restriction successively to the coordinate lines and planes shows that
the coefficients of $d\zeta_i$, $ds$, $ds\wedge d\zeta_i$, and
$ds\wedge d\zeta_1\wedge d\zeta_2$ vanish to order at least one in $t$.
Indeed, on every unimodular monomial chart the triple
$(d\log x_1,d\log x_2,d\log t)$ is a constant invertible linear combination
of the $d\log w_j$.  A coefficient $F(t)$ multiplying a nonzero logarithmic
residue is holomorphic only if its constant term vanishes: restricting to
$w_j=0$ gives that residue as $F(0)$, so $F(t)=tG(t)$.  Invertibility of the
monomial change of basis applies this argument to each displayed
coefficient.
Thus every coefficient used below descends with the asserted regularity at
all three points of the base.

The linear part of a deck transformation has the form
\[
 R_g=\begin{pmatrix}r_g&0\\ *&1\end{pmatrix},
 \qquad \jmath_g:=r_g^{-1}=c_g\tau+d_g,
 \qquad g^*(d\tau)=\jmath_g^{-2}d\tau.
\]
For a one-form, the second component $c_2$ is a holomorphic invariant
function on the base and vanishes at the cusp, hence is zero.  The second
column of $c\dot Z=0$ then gives $c_1\dot\tau=0$, so $c_1=0$.  What remains
is a holomorphic one-form on $\mathbb P^1$, and hence vanishes.

For a two-form, $b_2\,dz$ descends to a holomorphic one-form on
$\mathbb P^1$ and is zero.  The remaining form
$\psi=b_1\,dz$ satisfies $g^*\psi=\jmath_g\psi$.  Consequently
\[
 \psi^{\otimes2}\otimes d\tau
\]
is invariant.  At the order-three and order-four points it descends as a
cubic differential with a pole of order at most two, while the toric cusp
estimate gives a pole of order at most one.  It is therefore a section of
\[
 (\Omega^1_{\mathbb P^1})^{\otimes3}
 (2p_1+2p_2+p_0)\simeq\mathcal O_{\mathbb P^1}(-1),
\]
so it is zero.  Since $\dot\tau$ is nonzero on a dense open set,
$\psi=0$.  The coefficient of a three-form obeys the same weight-one law
and the same argument applies.  Hence all three spaces vanish.
\end{proof}
